\section{A measure on the continuum gives separable quotients}
\label{pos:section}\label{sec:positive}

We prove the positive implication over the real field.  The argument has two
steps. First we show that failure of a separable infinite-dimensional
quotient forces an evaluation image of maximal density, and use that image
to construct a dual family whose evaluation at each fixed vector is almost
surely zero. We then code the family by sequences in a separable space and
select an infinite sequence satisfying all finite suppression estimates.

\subsection{The measure hypothesis}

Throughout this section let $\kappa=\mathfrak c=2^{\aleph_0}$, regarded as its
initial ordinal.  Assume that $\mu$ is a probability measure on
$\mathcal P(\kappa)$, vanishes on singletons, and is \emph{$\kappa$-additive}
in the following convention: for every disjoint family $(A_\xi)_{\xi<\lambda}$
with $\lambda<\kappa$,
\[
 \mu\Bigl(\bigcup_{\xi<\lambda}A_\xi\Bigr)
   =\sum_{\xi<\lambda}\mu(A_\xi),
\]
where a sum of nonnegative terms is the supremum of its finite subsums.
In particular, $\mu$ is countably additive, and a union of fewer than
$\kappa$ null sets is null: disjointify such a family before applying the
displayed identity.  No additivity assertion for families of size
$\kappa$ is intended.

\begin{lemma}\label{pos:measure}
The cardinal $\kappa$ is regular.  The measure $\mu$ is atomless, and for
each positive integer $N$ the set $\kappa$ has a partition into $N$ sets
of measure $1/N$.
\end{lemma}

\begin{proof}
Every set of cardinality less than $\kappa$ is null.  If
$\operatorname{cf}(\kappa)<\kappa$, a cofinal family of fewer than $\kappa$
initial segments would therefore express $\kappa$ as a null set.

Fix an injection $b:\kappa\to\{0,1\}^{\N}$.  If $A$ were a positive atom,
then for each successive bit one of the two cells would have full measure
in the preceding cell of $A$.  This gives decreasing subsets $A_n\subset A$
of measure $\mu(A)$ on which the first $n$ bits are fixed.  Their intersection
has at most one point, whereas countable additivity gives it measure
$\mu(A)>0$.  This contradiction proves atomlessness.

We include the elementary divisibility argument.  Every positive set has a
positive subset of measure at most any given $\varepsilon>0$: repeatedly
split a positive set into two positive sets and retain the smaller part.
A maximal disjoint family of such small positive subsets of a measurable
set $A$ is countable, since the measure is finite, and covers $A$ modulo a
null set.  For $0\le t\le\mu(A)$, successive unions of these pieces either
attain measure $t$, converge upwards to $t$, or first cross $t$.  In the
last case the union before that crossing has measure in
$(t-\varepsilon,t)$.  Thus one can approximate $t$ from below with error at
most $\varepsilon$.  Apply this construction repeatedly in the remainder
of $A$, with errors tending to zero, and take the union.  It has measure
exactly $t$.  Removing sets of measure $1/N$ one at a time gives the
required partition.
\end{proof}

\subsection{The density forced by failure of a quotient}

The evaluation operator $R_E:X\to E^*$ for $E\subset X^*$ is as in
Proposition~\ref{ev:criterion}.  The following argument uses the
finite-dimensional Euclidean sections supplied by Dvoretzky's
Theorem~\cite{Dvoretzky1961,Figiel1976}.

\begin{lemma}\label{pos:density}
Let $X$ be a real Banach space with no separable infinite-dimensional
quotient, and let $E\subset X^*$ be an infinite-dimensional separable
closed subspace.  Under the measure hypothesis above, the set
$D=R_E(B_X)\subset B_{E^*}$ has norm density $\kappa$.
\end{lemma}

\begin{proof}
Every member of $E^*$ is determined by its values on a fixed countable
dense subset of $E$, so $|D|\le\mathfrak c=\kappa$.  Suppose that $A\subset D$
is norm dense and $|A|<\kappa$.

For $n\ge1$ put $N_n=16^n$.  Inductively choose vectors
$(u_{n,j})_{j=1}^{N_n}\subset E$ such that
\begin{equation}\label{pos:rows}
 \|a\|_2\le\Bigl\|\sum_{j=1}^{N_n}a_j u_{n,j}\Bigr\|
       \le2\|a\|_2\qquad(a\in\R^{N_n}).
\end{equation}
After choosing row $n$, choose a finite set $D_n\subset B_X$ that
$1/2$-norms the span of all rows up to $n$.  Require every subsequent row
to annihilate $D_n$.  At each step the finitely many annihilation conditions
leave an infinite-dimensional closed subspace of $E$, so Dvoretzky's
Theorem gives \eqref{pos:rows} there.

For each $n$, Lemma~\ref{pos:measure} gives a map
$j_n:\kappa\to\{1,\ldots,N_n\}$ whose fibers have measure $1/N_n$.
For $y\in A$, \eqref{pos:rows} implies
$\|(y(u_{n,j}))_{j=1}^{N_n}\|_2\le2$, and hence
\[
 \int_\kappa |y(u_{n,j_n(\alpha)})|\,d\mu(\alpha)
 \le \frac2{\sqrt{N_n}}.
\]
The sum of these expectations is finite.  Monotone convergence shows that
$\sum_n|y(u_{n,j_n(\alpha)})|<\infty$ for almost every $\alpha$.
Intersecting these conull sets over $y\in A$ is valid because $|A|<\kappa$.
Choose one $\alpha$ in their intersection and set $g_n=u_{n,j_n(\alpha)}$.

The vectors $g_n$ have norms between $1$ and $2$, and the sets $D_n$ are
actual half-normers of their prefix spans, annihilated by all later vectors.
Lemma~\ref{ev:prefix} therefore makes $(g_n)$ basic, with coefficient
functionals $g_n^\#\in G^*$ of norm at most $4$, where
$G=[g_n:n\ge1]$.  Set $Z=[g_n^\#:n\ge1]\subset G^*$.
For $y\in A$, the series
$\sum_n y(g_n)g_n^\#$ converges in $G^*$, since its tail has norm at most
$4\sum_{n>m}|y(g_n)|$.  On the dense span of the $g_n$ its sum agrees with
$y|_G$, so $y|_G\in Z$.  Restriction is contractive and $A$ is dense in
$D$, whence $R_G(B_X)\subset Z$.  Scaling gives $R_G(X)\subset Z$.
Lemma~\ref{ev:prefix} now says that $R_G:X\to Z$ is onto, contradicting
the hypothesis on $X$.
\end{proof}

Uniform sampling in this proof requires only the indicated marginal
distributions; the maps $j_n$ need not be independent.

The maximal density lets us separate each initial span in a transfinite
enumeration of $D$ from a point outside that span. The resulting dual
functional annihilates every earlier point of the enumeration. Thus, at
any fixed point, the exceptional indices form a set of cardinality less
than $\kappa$ and have measure zero.

\begin{lemma}\label{pos:family}
Let $X$ and $E$ satisfy the hypotheses of Lemma~\ref{pos:density}.
There are $\varepsilon>0$ and a family $(h_\alpha)_{\alpha<\kappa}$ in
$S_{E^{**}}$ such that the functionals
$f_\alpha=h_\alpha\circ R_E\in X^*$ satisfy
\begin{equation}\label{pos:seminormalized}
 \varepsilon<\|f_\alpha\|\le1\qquad(\alpha<\kappa),
\end{equation}
and, for every $y\in R_E(B_X)$,
\begin{equation}\label{pos:pointwise-null}
 \mu\{\alpha:h_\alpha(y)\ne0\}=0.
\end{equation}
\end{lemma}

\begin{proof}
Write $D=R_E(B_X)$.  There is a $\delta$-separated subset of $D$ of
cardinality $\kappa$ for some $\delta>0$.  Otherwise, choose a maximal
$1/n$-separated subset for each $n\ge1$.  Each would have cardinality
less than $\kappa$, and their union would be dense in $D$ with cardinality
less than $\kappa$, by regularity.  This contradicts
Lemma~\ref{pos:density}.

Enumerate $D=(y_\beta)_{\beta<\kappa}$.  For $\alpha<\kappa$, let
$M_\alpha=[y_\beta:\beta<\alpha]\subset E^*$.  Its density is less than
$\kappa$.  Some $y\in D$ has distance greater than $\delta/4$ from
$M_\alpha$: if every point of the separated set had distance at most
$\delta/4$, a dense subset of $M_\alpha$ of cardinality less than
$\kappa$ would provide centers at distance less than $3\delta/8$ from
all those points.  Two points cannot share such a center, a contradiction.
Hahn--Banach gives $h_\alpha\in S_{E^{**}}$ annihilating $M_\alpha$ and
satisfying $|h_\alpha(y)|>\delta/4$.  Since $y=R_Ex$ for some $x\in B_X$,
the functional $f_\alpha=h_\alpha R_E$ satisfies
\eqref{pos:seminormalized} with $\varepsilon=\delta/4$.
For $y=y_\beta$, its nonzero evaluations occur only at indices
$\alpha\le\beta$, a set of cardinality less than $\kappa$ and therefore
$\mu$-null.
\end{proof}

\subsection{Borel coding and ordered sampling}

We next explain how to use \eqref{pos:pointwise-null} when the vector
being evaluated is itself chosen using other indices.  This is the only
place where an interchange of integrals is needed.
Borel evaluation also plays a central role in the perfect-unconditionality
theorem of Argyros, Dodos and Kanellopoulos
\cite[Theorem~4]{ArgyrosDodosKanellopoulos2008}. Here it permits an
ordered sampling argument using ordinary Fubini. The next lemma needs
only a countably additive measure; the stronger additivity was used to
produce the family to which it will be applied.

\begin{lemma}\label{pos:sampling}
Let $X$ be a real Banach space, let $E\subset X^*$ be a separable closed
subspace containing no isomorphic copy of $\ell_1$, and put
$D=R_E(B_X)$.  Let $I$ carry a countably additive probability $\mu$ on
$\mathcal P(I)$ that vanishes on singletons.  Suppose that
$(h_\alpha)_{\alpha\in I}\subset E^{**}$ satisfies
$\mu\{\alpha\in I:h_\alpha(y)\ne0\}=0$ for every $y\in D$, and that
$f_\alpha=h_\alpha R_E$ satisfies
$0<\varepsilon\le\|f_\alpha\|\le M<\infty$ for every $\alpha\in I$.
Then $X^*$ contains an infinite unconditional basic sequence.
\end{lemma}

\begin{proof}
By the Odell--Rosenthal Theorem~\cite{OdellRosenthal1975}, each
$h_\alpha$ is the weak-star limit of a sequence from $E$.  Choose one
such sequence $c_\alpha\in E^{\N}$.  Equip $E^{\N}$ with the product of
the norm topologies and $K=B_{E^*}$ with the weak-star topology; both are
Polish spaces, and $K$ is compact.
There is a Borel function
\[
 \Phi:E^{\N}\times K\longrightarrow\R
\]
whose value at $(c,y)$ is $\lim_n y(c_n)$ if this limit exists and is
zero otherwise.  Indeed each map $(c,y)\mapsto y(c_n)$ is continuous,
and the Cauchy criterion makes its convergence set Borel.  To check the
joint continuity, use
$|y'(e')-y(e)|\le\|e'-e\|+|y'(e)-y(e)|$ for $y,y'\in K$.

The map $\alpha\mapsto c_\alpha$ is measurable because its domain has
the full power-set sigma algebra.  Let $\eta$ be its Borel pushforward
probability on $E^{\N}$.  The function
\begin{equation}\label{pos:H}
 H(y)=\int_{E^{\N}}\mathbf1_{\{\Phi(c,y)\ne0\}}\,d\eta(c)
\end{equation}
is Borel on $K$, by parameter integration of a bounded Borel function.
Our chosen codes give $\Phi(c_\alpha,y)=h_\alpha(y)$, so
$H(y)=0$ for every $y\in D$.

For $n\ge1$ define a probability $\lambda_n$ on all subsets of $I^n$
by ordered iteration, with the first variable outermost:
\[
 \lambda_n(A)=\int_I\cdots\int_I
      \mathbf1_A(\alpha_1,\ldots,\alpha_n)
      \,d\mu(\alpha_n)\cdots d\mu(\alpha_1).
\]
Put $\lambda_0$ equal to the unit mass on the one empty tuple.
Every section and every intermediate integrand is measurable.  Induction
and monotone convergence prove countable additivity of $\lambda_n$;
approximation by simple functions proves the same ordered integral formula
for bounded functions.  This construction does not assert that changing
the order of these power-set integrals preserves their value.

Fix $n$, rational coefficients $a_1,\ldots,a_n$, and a subset
$s\subset\{1,\ldots,n\}$.  For each tuple $(\alpha_j)_{j\in s}$ choose
$x_s\in B_X$ such that
\begin{equation}\label{pos:normer}
 \left|\sum_{j\in s}a_j f_{\alpha_j}(x_s)\right|
 \ge\frac12\left\|\sum_{j\in s}a_j f_{\alpha_j}\right\|.
\end{equation}
Choose $x_s=0$ if that sum is zero.  The choice depends only on the
indices in $s$.
For $i\notin s$ we claim
\begin{equation}\label{pos:cross-vanishing}
 f_{\alpha_i}(x_s)=0\quad\text{for $\lambda_n$-almost every tuple.}
\end{equation}

Fix any earlier prefix $(\alpha_1,\ldots,\alpha_{i-1})$.  The map
\[
 Y:I^{n-i}\longrightarrow K,\qquad
 Y(\alpha_{i+1},\ldots,\alpha_n)=R_Ex_s
\]
is independent of $\alpha_i$ and takes its values in $D$.  It is measurable
on the full power-set domain.  Let $\nu=Y_*\lambda_{n-i}$, with the same
interpretation when $i=n$.  The ordered integral, starting at variable
$i$, of the indicator of a nonzero evaluation in
\eqref{pos:cross-vanishing} is
\begin{align*}
 &\int_I\int_{I^{n-i}}
    \mathbf1_{\{\Phi(c_{\alpha_i},Y(t))\ne0\}}
           \,d\lambda_{n-i}(t)\,d\mu(\alpha_i)\\
 &\qquad=\int_{E^{\N}}\int_K
    \mathbf1_{\{\Phi(c,y)\ne0\}}\,d\nu(y)\,d\eta(c)
   =\int_K H(y)\,d\nu(y)=0.
\end{align*}
The change of order in the last line is ordinary Fubini for a Borel
function on the product of two Polish spaces with Borel probability
measures.  The final equality does not require $D$ to be Borel:
the Borel set $\{H\ne0\}$ has empty inverse image under $Y$, so is
$\nu$-null.  Since the earlier prefix was arbitrary, integrating over
those variables proves \eqref{pos:cross-vanishing}.

Combining \eqref{pos:normer} and \eqref{pos:cross-vanishing} gives
\begin{equation}\label{pos:suppression}
 \left\|\sum_{j\in s}a_j f_{\alpha_j}\right\|
 \le2\left\|\sum_{j=1}^n a_j f_{\alpha_j}\right\|
\end{equation}
almost surely.  Also the coordinates are pairwise distinct almost surely:
for $r<t$, the event $\alpha_t=\alpha_r$ constrains variable $t$ to a
singleton after its preceding variables have been fixed.  There are
countably many rational coefficient tuples and finitely many subsets $s$.
Consequently, for each $n$ there is a set $G_n\subset I^n$ of
$\lambda_n$-measure one on which distinctness and all these inequalities
hold.

We must choose one sequence whose prefixes all lie in these sets.
Suppose a prefix $p=(\alpha_1,\ldots,\alpha_m)$ has been chosen with
all its nonempty prefixes in their respective $G_j$, and with
\begin{equation}\label{pos:prefix-invariant}
 \lambda_{n-m}((G_n)_p)=1\qquad\text{for every }n>m.
\end{equation}
Here $(G_n)_p$ denotes the section obtained by fixing the prefix $p$.
The invariant holds for the empty prefix.  For $n=m+1$, it says that
$(p,\alpha)\in G_{m+1}$ for almost every $\alpha$.
For every $n>m+1$, the definition of the ordered measures gives
\[
 1=\lambda_{n-m}((G_n)_p)
  =\int_I\lambda_{n-m-1}((G_n)_{(p,\alpha)})\,d\mu(\alpha).
\]
The integrand lies in $[0,1]$, so it equals $1$ for almost every $\alpha$.
The intersection of these countably many conull sets is nonempty.
Choose the next index there.  This constructs the required infinite
sequence without an infinite product measure.

Continuity extends \eqref{pos:suppression} to real coefficients on every
finite prefix.  Singleton suppressions and $\|f_{\alpha_j}\|\ge\varepsilon$
give linear independence.  The initial-segment projections have norm
at most $2$ on the algebraic span and extend to its closure.  Approximation
by finite sums shows that these projections converge to the identity,
so the selected sequence is a Schauder basis for its closed span.
Every finite coordinate suppression also has norm at most $2$.
For a vector approximated within $\delta$ by a finite sum, every finite
coordinate sum outside that sum's support has norm at most $2\delta$.
It follows that the basis expansion converges unconditionally.
\end{proof}

\subsection{The positive conclusion}

\begin{theorem}\label{pos:theorem}
Suppose that $\mathfrak c$ carries a $\mathfrak c$-additive probability
measure on its full power set that vanishes on singletons, with additivity
understood for disjoint families of cardinality strictly less than
$\mathfrak c$.  Then every infinite-dimensional real Banach space has a
separable infinite-dimensional quotient.
\end{theorem}

\begin{proof}
Suppose that $X$ is a counterexample.  A consequence of the
Hagler--Johnson Theorem, in the precise form of
\cite[Proposition~16]{ArgyrosDodosKanellopoulos2008}, says that an
unconditional basic sequence in $X^*$ gives a separable
infinite-dimensional quotient of $X$.  Thus $X^*$ contains no such
sequence.  Choose an infinite-dimensional separable closed subspace
$E\subset X^*$.  It contains no isomorphic copy of $\ell_1$, whose usual
basis would give an unconditional basic sequence in $X^*$.

Lemmas~\ref{pos:density} and~\ref{pos:family} give a seminormalized family
$f_\alpha=h_\alpha R_E$ with the pointwise null-support property.
Lemma~\ref{pos:sampling} then produces an unconditional basic sequence in
$X^*$, the desired contradiction.
\end{proof}

The real positive implication is now complete. The next section transfers
it to complex Banach spaces; Section~\ref{model:section} supplies a model
of its measure hypothesis from a measurable cardinal.
